\documentclass[10pt,a4paper]{amsart}
\usepackage[margin=19mm]{geometry}
\usepackage{amsmath,amssymb,mathtools}
\usepackage[scr=rsfs,cal=euler]{mathalfa}
\usepackage{xfrac}
\usepackage[unicode,hidelinks,bookmarks=false]{hyperref}
\newcommand{\ie}{i.e.\ }
\newcommand{\cf}{cf.\ }
\newcommand{\resp}{respectively\ }
\newcommand{\Subsection}{\subsection}
\providecommand{\nobreakdash}{\nobreak\hbox{-}}
\setlength{\emergencystretch}{2em}
\multlinegap0pt
\usepackage{amssymb}
\usepackage{enumerate}
\newtheorem{theorem}{Theorem}
\title{Mixing Times and Cutoff for the Rook's Walk}
\author{Jonatan Kaare-Rasmussen}
\date{Source: arXiv:2604.07478v1 (CC BY 4.0)}
\begin{document}
\maketitle

\noindent\emph{Source and licence.} Mixing Times and Cutoff for the Rook's Walk. Version arXiv:2604.07478v1 (CC BY 4.0), distributed by arXiv under the Creative Commons Attribution 4.0 International licence, http://creativecommons.org/licenses/by/4.0/, as recorded in the arXiv licence metadata for this exact version and shown on the abstract page https://arxiv.org/abs/2604.07478v1. Full licence notice: \texttt{licenses/arxiv-2604.07478-CC-BY-4.0.txt}.\par\smallskip

\noindent\textit{Editorial scope.} Counted unit: the article's theorem main1 (source0.tex lines 65--72). The article's complete original proof of it is delivered with it (source0.tex lines 426--443, one \texttt{proof} environment of 18 lines, which the article places away from the statement). No mathematics is written, repaired or re-derived: the statement and the proof are the article's own lines, in the article's own notation, and no retained line is substituted or re-flowed.  Retained uncounted as the source-local context the proof needs: lines 389--395.  Nothing was omitted from any retained span, so the package carries no omission record.  No retained line contains a citation, so the wrapper carries no bibliography and no bibliography entry.  No retained line contains a figure environment, an image-inclusion command or drawing code, so no image is delivered and the package declares no asset dependency.  Editorial formatting only, in addition: an AMS \texttt{amsart} preamble that restates the article's own declarations and macros used by the retained lines, namely the environments \texttt{theorem} on separate counters, together with the standard packages those lines need.  The retained environments print in this document as Theorem 1 in order of appearance, whereas the article prints the same items with the numbering of its own full document.  The complete native source of the article as distributed in the arXiv e-print for this exact version is bundled unchanged as \texttt{sources/198093/source0.tex}.
\begin{theorem}
   \label{wilson}
   Let $(X_t)$ be an ergodic Markov chain with state space $\Omega$ and transition matrix $P$. Let $\Phi$ be an eigenfunction of $P$ with real eigenvalue $\alpha$ satisfying $1/2 < \alpha< 1$. Let $R>0$ such that:
   \[\mathbb{E}_x\left(|\Phi(X_1)-\Phi(x)|^2\right)\leq R,\quad \forall x\in \Omega\]
   Then for any $x\in \Omega$ and $\epsilon\in (0,1)$,
   \[t_\text{mix}(\epsilon) \geq \frac{1}{2\log(1/\alpha)}\left[\log \left(\frac{(1-\alpha)\Phi(x)^2}{2R}\right) + \log\left(\frac{1-\epsilon}{\epsilon}\right)\right]\]
\end{theorem}

      \begin{theorem}
         \label{main1}
      Let $t_\text{mix}(\epsilon)$ with $\epsilon >0$ be the mixing time of the Rook's Walk with dimension $d$ and board-length $n$. Then,
      \begin{multline*}
       \left(\log \left(\frac{d(n-1)}{2n}\right) + \log\left(\frac{1-\epsilon}{\epsilon}\right)\right)\left(2\log\left(\frac{d(n-1)}{d(n-1)-n}\right)\right)^{-1} \\
        \leq t_\text{mix}(\epsilon)  \leq-\frac{d(n-1)}{2n}\log\left((4\epsilon^2+1)^{1/d}-1\right) 
      \end{multline*}
   \end{theorem}

   \begin{proof}[Proof of Lower Bound in Theorem \ref{main1}]
      From Theorem \ref{wilson} we have,
   \[t_\text{mix}(\epsilon) \geq \frac{1}{2\log(1/\alpha)}\left[\log \left(\frac{(1-\alpha)\Phi(x)^2}{2R}\right) + \log\left(\frac{1-\epsilon}{\epsilon}\right)\right], \quad \forall x\in \Omega\]
   where $\alpha = 1-\frac{n}{d(n-1)}$, $R=\frac{n^2}{d^2(n-1)^2}$ and $\Phi(x) = 1-(1-\alpha)x$. Maximizing over $x\in \Omega$,
   \begin{align*}
   t_\text{mix}(\epsilon) &\geq \frac{1}{2\log(1/\alpha)}\left[\log \left(\frac{(1-\alpha)\Phi(0)^2}{2R}\right) + \log\left(\frac{1-\epsilon}{\epsilon}\right)\right]\\
   &= \frac{1}{2\log(1/\alpha)}\left[\log \left(\frac{1-\alpha}{2R}\right) + \log\left(\frac{1-\epsilon}{\epsilon}\right)\right]
   \end{align*}
   Consider just $\frac{1-\alpha}{2R}$,
   \begin{align*}
      \frac{1-\alpha}{2R}&= \frac{\frac{n}{d(n-1)}}{2\frac{n^2}{d^2(n-1)^2}}\\
      &= \frac{nd^2(n-1)^2}{2d(n-1)n^2}\\
      &= \frac{d(n-1)}{2n}
   \end{align*}
   Hence,
   \[t_\text{mix}(\epsilon) \geq \left(\log \left(\frac{d(n-1)}{2n}\right) + \log\left(\frac{1-\epsilon}{\epsilon}\right)\right)\left(2\log\left(\frac{d(n-1)}{d(n-1)-n}\right)\right)^{-1}\]
   as required.
   \end{proof}

\end{document}
